\section{The Availability Gap}\label{sec:core}

\subsection{The placebo comparison}\label{sec:core-placebo}

To separate withheld from non-existent, we require a control group of non-endowment \citep{dye1985}. Dates appear to provide one: a postdating benchmark could not be reported without prepublication access, whereas an eligible one could. A disclosure label would report whether it was, and none enters. Any innocent reason an eligible looks poor, difficulty drift, saturation, or measure change, operates equally well on both groups, meaning their difference should be zero. One such innocent reason is inherently asymmetric: a benchmark prior to shipping can lie in the training set for a given model and one after it cannot, so contamination increases eligible standing only, beyond any reweighting of the comparison.

Not zero. On average eligible beats postdating by 12.23 percentile points among 146 releases of both cells, with the median 10.69, and the proportion of releases where it wins is 78.1\%. This computation does not know whether benchmarks are reported. By constructing one that treats both sides of the window equally (see below), we find eligible to beat by 9.14 percentile points across 141 releases, positive 70.9\%, still without removing it. Using a cluster sign-flip randomization test of the difference at the 18 organisations with both cell types produces $p = 0.0001$ as the smallest available from the draws. Groups ineligible to be reported already stand lower than ones eligible, so it is not a control.

\subsection{Decomposing the gap}\label{sec:core-boundary}

In \autoref{tab:decomposition} we identify the failure using a placebo regression with an indicator for whether the benchmark postdates the release (always within release). With release effects alone this is $-13.5$ points. But benchmark fixed effects account for 63.4\% of it, structurally since the placebo cells are constructed to exist primarily in those benchmarks arriving late, which are less saturated, with 8.3\% of scores in the latter half of the 54 benchmarks with $\geq 10$ scores in the top decile of its own range vs.\ 13.2\% in the former ($p = 0.015$), leaving more room below the ceiling. The term capturing the asymmetry in peer-window, i.e.\ what share of comparisons are made to models released after the model of interest, controls for 39.8\% on its own, and once we include both benchmark fixed effects and both window terms, the remaining is 0.003 points, with a confidence interval up to $+3.4$, which is no more than one third of the pre-control difference.

\begin{table}[t]
\centering
\small
\begin{tabular}{lrrr}
\toprule
Conditioning set & Placebo coefficient & $t$ & Absorbed \\
\midrule
Release fixed effects only & $-13.443$ (1.238) & $-10.855$ & 0.0\% \\
\quad + benchmark fixed effects (two-way absorbed) & $-4.869$ (1.260) & $-3.865$ & 63.8\% \\
\quad + peer-window asymmetry & $-9.515$ (1.254) & $-7.587$ & 29.2\% \\
\quad + peer count & $-2.566$ (1.305) & $-1.967$ & 80.9\% \\
Two-way fixed effects + both window terms & $+0.246$ (1.618) & $+0.152$ & 98.2\% \\
\bottomrule
\end{tabular}
\caption{%
Each row is the coefficient on a placebo indicator in a regression of within-benchmark standing on that indicator, always within release, so the release's own capability level is differenced out throughout. Provider-clustered standard errors in parentheses over 46 organisations; $n = 2831$ cells (2355 eligible, 476 placebo). Absorbed is the reduction in the absolute coefficient relative to the first row. Peer count is a suppressor when entered alone, so the final row is joint conditioning on both window terms and must not be read as peer count explaining a remainder.}
\label{tab:decomposition}
\end{table}



For the second dimension, the mechanism is geometric, with the peer-window being symmetric in calendar time, 182 days before or after, but not model coverage since the coverage begins when the benchmark arrives. \autoref{fig:boundary} depicts both objects against benchmark maturity at release: the relative proportion of newer peers declines with maturity and jumps at the availability threshold, whereas standing runs through it without a detectable break. Thus, the mechanism is discontinuous where the outcome is not; composition is a property of timing alone.

The share of a cell's window released after the focal model is 0.4683 for eligible cells and 0.7065 for placebo cells, against 0.5 for a two-sided window. Within release, the slope of standing on that share is $-30.63$ percentile points per unit among eligible cells alone. \autoref{fig:one-curve} plots standing against that share for both groups at once. Eligible and placebo cells do not form two clouds; they trace a single declining curve, common only once benchmark composition is absorbed (\autoref{tab:decomposition}), and differ in where along it they fall. On the variable that assigns them, benchmark maturity at release, they cannot overlap, so the contrast is an extrapolation across that boundary rather than a comparison within common support. The check is at the boundary, the one place Section \ref{sec:data-model} leaves identified under continuity \citep{hahn2001identification}. Nothing detectable happens there. Local linear estimates are negative at every bandwidth and order, every interval covers zero though the widest reaches $+8.05$, and a window chosen on covariate balance puts the difference at $-2.52$, randomization $p$ 0.447 \citep{cattaneo2015randomization}. The two groups' mean maturity differs by 793 days; the gap lives in that distance, not at the boundary (balance tests in Appendix~\ref{sec:supporting}).

Standing decomposes into a weighted average of position among older and newer peers, $\text{percentile} = (1 - w)\,P_{\text{old}} + w\,P_{\text{new}}$ with $w$ the newer share, exact cell by cell. Replacing the empirical weight with one half defines the side-balanced measure, available for 92.5\% of all cells and 92.8\% of eligible-or-placebo cells, and cuts the same slope to $-8.56$.

That the identity holds does not make the gradient an accounting tautology. Shuffling scores within benchmark destroys the capability trend while leaving dates, window membership, peer counts and newer shares bit-identical. Under that null, the slope has mean $-0.09$ and standard deviation 3.12. The observed $-30.63$ sits 9.8 standard deviations away, and none of 300 draws reaches it. Replacing scores instead with a within-benchmark linear fit on date plus permuted residuals, so that a capability trend exists and nothing else does, recovers 93\% of the slope over 200 draws. Comparison-set geometry is inert on its own and becomes a bias only when it multiplies a real trend in capability.

\begin{figure}[!htb]\centering
\begin{tikzpicture}
  \begin{axis}[causalre wide, height=2.8cm, xmin=-360, xmax=720, xtick={-360,0,360,720}, xticklabels={,,,}, ylabel={newer share}, ymin=0.25, ymax=0.95, ytick={0.3,0.5,0.7,0.9}, name=mech]
    \addplot[black, dashed, thin, forget plot] coordinates {(0,0.25) (0,0.95)};
    \addplot[cSF, only marks, mark=square*, mark size=1.4pt,
             error bars/.cd, y dir=both, y explicit]
      coordinates {(-270.000,0.82) +- (0,0.08) (-210.000,0.79) +- (0,0.08) (-150.000,0.73) +- (0,0.11) (-90.000,0.71) +- (0,0.08) (-30.000,0.63) +- (0,0.10)};
    \addplot[cNYC, only marks, mark=*, mark size=1.4pt,
             error bars/.cd, y dir=both, y explicit]
      coordinates {(30.000,0.60) +- (0,0.07) (90.000,0.53) +- (0,0.08) (150.000,0.50) +- (0,0.09) (210.000,0.42) +- (0,0.08) (270.000,0.41) +- (0,0.10) (330.000,0.43) +- (0,0.11) (390.000,0.41) +- (0,0.10) (450.000,0.48) +- (0,0.10) (510.000,0.44) +- (0,0.10) (570.000,0.31) +- (0,0.11) (630.000,0.40) +- (0,0.11) (690.000,0.37) +- (0,0.08)};
  \end{axis}
  \begin{axis}[causalre wide, height=2.8cm, xmin=-360, xmax=720, xtick={-360,0,360,720}, xlabel={benchmark maturity at release (days)}, ylabel={standing}, ymin=30, ymax=75, at={(mech.south west)}, anchor=north west, yshift=-12pt,
    legend to name=leg:boundary, legend columns=2]
    \addplot[black, dashed, thin, forget plot] coordinates {(0,30) (0,75)};
    \addplot[cSF, only marks, mark=square*, mark size=1.4pt,
             error bars/.cd, y dir=both, y explicit]
      coordinates {(-270.000,40.78) +- (0,18.37) (-210.000,41.55) +- (0,14.24) (-150.000,45.00) +- (0,16.04) (-90.000,41.36) +- (0,14.81) (-30.000,44.91) +- (0,13.75)};
    \addlegendentry{benchmark postdates the release}
    \addplot[cNYC, only marks, mark=*, mark size=1.4pt,
             error bars/.cd, y dir=both, y explicit]
      coordinates {(30.000,45.04) +- (0,13.16) (90.000,46.87) +- (0,13.81) (150.000,54.42) +- (0,11.65) (210.000,51.49) +- (0,11.25) (270.000,53.77) +- (0,12.19) (330.000,50.14) +- (0,13.47) (390.000,58.21) +- (0,12.71) (450.000,46.04) +- (0,12.76) (510.000,55.12) +- (0,13.60) (570.000,57.18) +- (0,13.73) (630.000,51.43) +- (0,12.65) (690.000,53.95) +- (0,12.64)};
    \addlegendentry{benchmark predates the release}
  \end{axis}
\end{tikzpicture}
\\[2pt]
{\footnotesize \ref{leg:boundary}}

\caption{The mechanism and the outcome at the availability boundary. Top: the share of a
cell's comparison window released after the focal model, against benchmark maturity at
release; it falls with maturity and jumps at the boundary. Bottom: within-benchmark
standing on the same axis, continuous through it. Sixty-day bins; intervals clustered on
organisation.}
\label{fig:boundary}
\end{figure}

\begin{figure}[!htb]\centering
\begin{tikzpicture}
  \begin{axis}[causalre body, xmin=0, xmax=1, xtick={0,0.5,1}, xlabel={share of window newer than the focal model}, ymin=20, ymax=80,
    ylabel={within-benchmark standing}, title={(a)},
    legend to name=leg:onecurve, legend columns=3]
    \addplot[cMUTE, line width=2.2pt, smooth, forget plot]
      coordinates {(0.050,60.90) (0.150,61.31) (0.250,57.85) (0.350,53.26) (0.450,52.33) (0.550,46.42) (0.650,44.05) (0.750,40.62) (0.850,41.55) (0.950,28.06)};
    \addplot[black, dashed, thin, forget plot] coordinates {(0.5,20) (0.5,80)};
    \addplot[cNYC, only marks, mark=*, mark size=1.4pt,
             error bars/.cd, y dir=both, y explicit]
      coordinates {(0.050,61.50) +- (0,12.43) (0.150,61.72) +- (0,13.22) (0.250,57.82) +- (0,10.51) (0.350,53.48) +- (0,10.97) (0.450,52.57) +- (0,9.92) (0.550,46.30) +- (0,9.40) (0.650,45.17) +- (0,10.66) (0.750,40.39) +- (0,12.27) (0.850,43.18) +- (0,13.62)};
    \addlegendentry{benchmark predates the release}
    \addplot[cSF, only marks, mark=square*, mark size=1.4pt,
             error bars/.cd, y dir=both, y explicit]
      coordinates {(0.350,49.27) +- (0,22.78) (0.450,49.40) +- (0,16.17) (0.550,46.96) +- (0,13.89) (0.650,40.71) +- (0,15.04) (0.750,40.99) +- (0,13.29) (0.850,40.43) +- (0,14.91)};
    \addlegendentry{benchmark postdates the release}
    \addplot[cMUTE, line width=2.2pt] coordinates {(0,0)};
    \addlegendentry{all cells pooled}
  \end{axis}
  \begin{axis}[causalre body, xmin=0, xmax=1, xtick={0,0.5,1}, xlabel={share of window newer than the focal model}, ymin=35, ymax=65, xshift=0.50\textwidth,
    title={(b)}, legend to name=leg:correction, legend columns=2]
    \addplot[cSF, smooth, mark=*, mark size=1.4pt]
      coordinates {(0.050,53.05) (0.150,55.63) (0.250,54.01) (0.350,53.06) (0.450,52.26) (0.550,48.46) (0.650,50.81) (0.750,45.60) (0.850,43.13)};
    \addlegendentry{windowed percentile}
    \addplot[cNYC, smooth, mark=square*, mark size=1.4pt]
      coordinates {(0.050,47.81) (0.150,52.05) (0.250,49.32) (0.350,49.54) (0.450,50.57) (0.550,47.87) (0.650,51.82) (0.750,46.51) (0.850,44.84)};
    \addlegendentry{side-balanced}
  \end{axis}
\end{tikzpicture}
\\[2pt]
{\footnotesize \ref{leg:onecurve} \hfill \ref{leg:correction}}

\caption{Standing against the share of a cell's comparison window released after the focal
model; dashed verticals mark each group's mean share, and the black dashed line the symmetric 0.5. Eligible and placebo cells lie on one
declining curve, so the contrast between them is a position on that curve rather than a fact
about disclosure. Intervals are clustered on organisation.}
\label{fig:one-curve}
\end{figure}

\subsection{Trimming and reweighting}\label{sec:core-recentring}

The phenomenon has a name in nonparametric estimation: local averages are biased near the edge of a support, and the standard remedies recentre rather than discard \citep{gasser1979,fan1992}. What is unusual is the boundary's coincidence with the endowment: a benchmark was available to a release exactly when the model did not sit at the left edge of that benchmark's coverage. The indicator of what a provider could have held is therefore the indicator of where the window is one-sided, and removing the edge removes the contrast. The coincidence has an older name, the age-period-cohort identity: release date, benchmark vintage and their difference are exactly collinear, and each has a measured signature here: $t_i$ in the capability trend, $\tau_b$ in saturation, $t_i - \tau_b$ in window composition. Only the linear components are unidentified in that structure \citep{fosse2019bounding,fosse2019review}, so what fails is a slope. Three sign restrictions bound it: capability rises with the calendar, later benchmarks are less saturated, and standing rises with maturity. All three are assumptions on unidentified components, each motivated by a measurement of an identified quantity. Under them the unidentified component accounts for at most 24 to 34\% of the placebo gap across specifications, roughly three percentile points, which is where the saturated regression also lands. The remainder is carried by the identified reduced-form slopes, and Appendix~\ref{sec:supporting} gives the derivation.

\begin{table}[H]
\centering
\small
\setlength{\tabcolsep}{4pt}
\begin{tabular}{lrrrrrr}
\toprule
Standing measure & Mean & Median & Positive & Sign-flip $p$ & Releases & Cells \\
\midrule
Windowed percentile, $\pm$182d & $+12.23$ & $+10.69$ & 78.1\% & 0.0001 & 146 & 100.0\% \\
Rank against all models ever & $+22.75$ & $+22.66$ & 89.0\% & 0.0001 & 146 & 100.0\% \\
Trim to two-sided windows & $+9.96$ & $+10.69$ & 72.2\% & 0.0001 & 115 & 70.4\% \\
Side-balanced percentile & $+9.14$ & $+9.62$ & 70.9\% & 0.0001 & 141 & 92.8\% \\
\bottomrule
\end{tabular}
\caption{%
The label-free placebo null under each candidate measure: release-level mean standing of eligible minus postdating benchmarks, in percentile points. Every entry should be zero. Sign-flip $p$ is the provider-clustered test of Section~\ref{sec:data-inference}, against zero; Cells, the share of cells carrying the measure.}
\label{tab:remedies}
\end{table}


\autoref{tab:remedies} shows what follows. Ranking against every model ever scored nearly doubles the contamination; trimming to two-sided windows discards 30\% of the cells to remove 19\%, its row the uncorrected gap on the retained cells, composition rather than correction. Shortening the window does reduce the gap, to $+8.59$ at 91 days, but in proportion to the asymmetry it removes rather than by breaking the mechanism. Across four widths and four benchmark-count minima the gap is positive in
all 16 specifications, from $+8.59$ to $+21.66$, and tracks the eligible-minus-placebo gap in window composition at $r = 0.991$. Dropping any one of the 18 organisations leaves it between $+11.45$ and $+13.20$. The correction is a covariate, not a subsample.

One reading of that number would be too generous to it. Window composition is a function of the same dates that define the placebo indicator, so conditioning on it does not adjust for a confounder sitting outside the comparison; it splits the gap into the channel carrying it. A study that holds disclosure labels and wants an estimate of $\theta$ must first argue that window composition is settled before a provider decides what to report, or it will condition away part of what it is measuring. Provider identity explains little of the variance in window composition; Appendix~\ref{sec:supporting} reports the decomposition and the coverage tilt it does not correct.

The side-balanced percentile is the one change of measure doing substantial work, cutting the placebo mean to $+9.14$ and the eligible-cell gradient to $-8.56$. \autoref{fig:correction} in Appendix~\ref{sec:supporting} puts it against the uncorrected measure, near flat where cells are dense. Reweighting cannot reach what remains: balancing repairs the weight on each side of a model but not which models an evaluator chose to score early, so evaluator attention stays a confound. Undefined wherever a side is empty, it accompanies rather than replaces the uncorrected measure.
